Fundstellen werden in einer Literaturdatenbank gesammelt. Das Wort
"`Datenbank"' mag einige falsche Vorstellungen wecken. Für \LaTeX\ versteht man
darunter eine schlichte Textdatei im sog. \hologo{BibTeX}-Format.
Listing~\ref{lst:latex:bibliography:bibtexentry} zeigt einen beispielhaften
Eintrag für dieses Handbuch.

\lstinputlisting[language=BibTeX,style=block,float,caption={Ein \hologo{BibTeX}-Eintrag},label={lst:latex:bibliography:bibtexentry}]{code/latex_bibliography_bibtexentry.bib}

Ein Eintrag besteht aus drei Teilen.
%
\begin{compactenum}
\item dem Typ, hier \texttt{@manual}
\item einem Schlüssel, hier \texttt{urz\_vorlage\_manual}
\item mehreren Attributen, die die Fundstellen näher beschreiben
\end{compactenum}

Natürlich kann man die Datei manuell mit einem Texteditor pflegen. Es gibt
allerdings auch spezialisierte Software, die die Bearbeitung einer
\hologo{BibTeX}-Datenbank erleichtert. Als Beispiel wären zu nennen:
%
\begin{compactitem}
\item K\hologo{BibTeX} (\url{https://userbase.kde.org/KBibTeX})
\item Zotero (\url{https://www.zotero.org})
\item JabRef (\url{https://www.jabref.org})
\end{compactitem}
%
Mitunter speichern diese Programme die Literaturdatenbank in einem eigenen
Format und exportieren die Fundstellen nach der Bearbeitlung lediglich im
\hologo{BibTeX}-Format.

In den seltesten Fälle ist es notwendig, eine \hologo{BibTeX}-Datenbank
händisch anzulegen. Verlage und Bibliotheken halten in ihren Katalogen meist
eine entsprechende Exportfunktion vor, die es ermöglicht, vorgefertigte
\hologo{BibTeX}-Einträge herunterzuladen. I.d.R. sind jedoch anschließende
Anpassungen notwendig, falls Umlaute oder Sonderzeichen in den Autorennamen
nicht korrekt kodiert sind oder Datenfelder fehlen.
